\section{Evaluation}\label{evaluation}

\begin{center}\fbox{\begin{minipage}{0.94\linewidth}\setlength{\parskip}{4pt}\small

\textbf{Gap: LEGO-Bench.} This section will report SCORE's scores on LEGO-Bench \citep{li2026lego}, a benchmark of scenes rebuilt from pictures. It waits until the benchmark's code is released and until SCORE's scene export and its Blender adapter are built.

\end{minipage}}\end{center}

\subsection{With and without the agent tools}\label{with-and-without-the-agent-tools}

The tools that check an agent's work (the completion gate, the resting triage with its drop test, the pick lock and the render-and-compare, Section 2) were tested on three places of world 1: a room (the workshop, 814 objects), an outdoor place (the street, 1,198 objects) and a dense room (the hangar, 1,657 objects). Each arm started from the same exported stage and the same inventory, and the same agent (Claude Opus 5.5 in Claude Code, headless) got the same prompt: make the place as complete as it can in about 50 minutes, through the place's layout data only. One arm had the framework as it was before these tools; the other had them. Both arms were measured afterwards with the same instruments (Table 2). Each arm ran once, so the table shows what happened, not an average.

\textbf{Table 2.} One session per arm and place. Faults are the resting triage's real faults left on the final stage, measured against the arm's own data; placeholders by the placeholder check; rows are inventory rows with no object and no written reason. Agent cost is the session's model cost in dollars; cloud is what its rented machines cost.

{\def\LTcaptype{none} % do not increment counter
\begin{longtable}[]{@{}
  >{\raggedright\arraybackslash}p{(\linewidth - 18\tabcolsep) * \real{0.1000}}
  >{\raggedright\arraybackslash}p{(\linewidth - 18\tabcolsep) * \real{0.1000}}
  >{\raggedright\arraybackslash}p{(\linewidth - 18\tabcolsep) * \real{0.1000}}
  >{\raggedright\arraybackslash}p{(\linewidth - 18\tabcolsep) * \real{0.1000}}
  >{\raggedright\arraybackslash}p{(\linewidth - 18\tabcolsep) * \real{0.1000}}
  >{\raggedright\arraybackslash}p{(\linewidth - 18\tabcolsep) * \real{0.1000}}
  >{\raggedright\arraybackslash}p{(\linewidth - 18\tabcolsep) * \real{0.1000}}
  >{\raggedright\arraybackslash}p{(\linewidth - 18\tabcolsep) * \real{0.1000}}
  >{\raggedright\arraybackslash}p{(\linewidth - 18\tabcolsep) * \real{0.1000}}
  >{\raggedright\arraybackslash}p{(\linewidth - 18\tabcolsep) * \real{0.1000}}@{}}
\toprule\noalign{}
\begin{minipage}[b]{\linewidth}\raggedright
Place, arm
\end{minipage} & \begin{minipage}[b]{\linewidth}\raggedright
Real faults (start → end)
\end{minipage} & \begin{minipage}[b]{\linewidth}\raggedright
Placeholders
\end{minipage} & \begin{minipage}[b]{\linewidth}\raggedright
Rows missing
\end{minipage} & \begin{minipage}[b]{\linewidth}\raggedright
Minutes
\end{minipage} & \begin{minipage}[b]{\linewidth}\raggedright
Agent cost
\end{minipage} & \begin{minipage}[b]{\linewidth}\raggedright
Output tokens
\end{minipage} & \begin{minipage}[b]{\linewidth}\raggedright
Cloud
\end{minipage} & \begin{minipage}[b]{\linewidth}\raggedright
Exports
\end{minipage} & \begin{minipage}[b]{\linewidth}\raggedright
Gate said unknown
\end{minipage} \\
\midrule\noalign{}
\endhead
\bottomrule\noalign{}
\endlastfoot
Workshop, without & 18 → 13 & 0 → 0 & 1 → 0 & 36 & \$2.26 & 30.4k & €0.85 & 7 & not used \\
Workshop, with & 18 → 0 & 0 → 0 & 1 → 0 & 26 & \$1.42 & 14.3k & €0.85 & 2 & 2 of 4 calls \\
Street, without & 11 → 0 & 0 → 0 & 0 → 0 & 38 & \$2.40 & 30.0k & €0.32 & 9 & not used \\
Street, with & 12 → 0 & 0 → 0 & 0 → 0 & 39 & \$1.21 & 10.4k & €0.76 & 1 & 0 of 5 calls \\
Hangar, without & 14 → 4 & 5 → 0 & 0 → 0 & 47 & \$2.01 & 26.6k & €0.00 & 7 & not used \\
Hangar, with & 14 → 9 & 5 → 0 & 0 → 0 & 32 & \$2.34 & 29.1k & €0.64 & 5 & 3 of 6 calls \\
\end{longtable}
}

The tools made building cheaper more often than they made the scene better. On the workshop and the street, the arm with the tools used about half the agent's output and cost, and exported the stage once or twice instead of seven to nine times, because the triage named the few real faults among hundreds of alarms. On the workshop it left no real fault, but 13 of its 18 were cleared by writing what holds a piece up into the inventory (the tools on the tool board hang from it; the arm monitors are clamped to their bench), which the creator still has to confirm; the other five were moved off the wall they stood in. On the hangar the arm without the tools left fewer faults: the arm with them spent its time replacing five placeholder slabs with the kit's own wall panels and ribs and trying the drop test, which failed on a room with no ground mesh (fixed after the run), and the nine faults it left are models with stray parts below their feet that no pose can fix. In no arm did the completion gate's hook have to refuse a claim, because no agent claimed a place was complete; both arms reported what was left. Every unknown answer of the gate was a check not yet run on the stage as it then was, which is what unknown means; the agents ran the checks and went on. Concept fit and the per-part surface check are not in the table, because their tools are not built yet, and no row carries an owner's pick, so the pick lock was never tested by an agent.

\begin{center}\fbox{\begin{minipage}{0.94\linewidth}\setlength{\parskip}{4pt}\small

\textbf{Gap: ablations.} This section will compare the framework with and without the optional world step, and with the checks before spending switched off, on the same places (the agent tools are compared above). It will measure how much of the concept is covered, how dense the built place is, the cost, the faults found after spending, and the creator's review. It waits until every stage runs from the framework repository.

\end{minipage}}\end{center}

\subsection{Style coherence}\label{style-coherence}

Whether a place keeps one style is read from the review pages' renders of its generated models, which draw every model under the same light, against the same background and by the same camera rule, so that two renders differ only by the model. Three readings come from the pixels: how well the coloured pixels fit the place's palette, how much the models' lightness varies, and how blotchy their surfaces are, as the 75th percentile of the lightness spread within small tiles. A fourth comes from DINOv2 \citep{oquab2023dinov2}: how alike two models of one place are, against two models of different places, with the same pairs compared again as plain silhouettes and that difference subtracted, so that what an object is does not count as its style. Over world 1's 140 generated models in 16 places, DINOv2 finds models of one place more alike than models of different places, but almost all of that is shape: the silhouettes show nearly the same difference, and what is left for the look is 0.009 on the cosine scale, as a mean over models (95 percent bootstrap interval 0.003 to 0.015). Every place of world 1 draws on one surface library, so this reading shows how distinct the places are rather than how consistent each one is. The palette reading says almost nothing on renders: most models are greys, steels and dark panels, the places' palettes overlap almost entirely, and even the generator's own colours fit them.

The surface reading does respond to the method. On the 54 generated models of the garage, the hangar and the lab, each rendered once painted from the library and once with the colours of the picture it was generated from, the painted render was less blotchy for 43 of the 54 and more even from model to model, the difference that Figure 4b shows by eye. The two renders differ in geometry as well as paint, because the picture-coloured one is the generator's own mesh, so this does not isolate the paint (Appendix B). None of the readings has yet been checked against the creator's reviews: since its first room, the hub, which the creator accepted (Table C1), world 1 has not been reviewed, so there is nothing yet for a reading to separate, and the most weathered place, the old station, reads the most blotchy because its wear is intended. Until reviews of both kinds exist, these are measurements of renders, not a measure of style.

\begin{center}\fbox{\begin{minipage}{0.94\linewidth}\setlength{\parskip}{4pt}\small

\textbf{Gap: robots in the simulation world.} This section will show a rigged robot loaded together with the simulation world in a robotics simulator, with masses, friction and joints written into the scene, and with motion from the framework's motion models. It waits on world 3.

\end{minipage}}\end{center}
